\section{Il paradigma FaaS in ambienti cloud-edge}
\label{sec:faas_cloud_edge}

Negli ambienti cloud-edge l'elaborazione è distribuita su un insieme
eterogeneo di nodi di calcolo, caratterizzati da differenti capacità
computazionali, disponibilità di memoria, ampiezza di banda e latenza
di rete. Nel paradigma Function-as-a-Service, le applicazioni sono
scomposite in funzioni stateless invocate su richiesta, la cui
esecuzione è demandata alla piattaforma. In un contesto cloud-edge,
ciascuna funzione può essere eseguita localmente su un nodo edge,
oppure instradata verso altri nodi o verso il cloud, in funzione
dello stato delle risorse disponibili al momento dell'invocazione.

Le differenze strutturali tra nodi cloud e nodi edge rendono il
consumo energetico un fattore particolarmente rilevante in questi
scenari. I nodi edge dispongono tipicamente di risorse limitate e
operano in condizioni di maggiore variabilità del carico, spesso
in assenza di infrastrutture di raffreddamento dedicate. Di
conseguenza, le decisioni di allocazione e gestione delle funzioni
non possono basarsi esclusivamente su metriche prestazionali, ma
devono tenere conto dell'impatto energetico delle scelte effettuate,
soprattutto in presenza di vincoli operativi stringenti.

Le tecniche discusse nelle sezioni precedenti — scheduling
energeticamente consapevole, meccanismi applicativi basati su approximate
computing e approcci carbon-aware — sono state sviluppate prevalentemente
in contesti cloud o HPC, con assunzioni sull'infrastruttura che non
sempre si trasferiscono direttamente all'ambiente cloud-edge. In
particolare, la modularità della piattaforma serverless di riferimento
è un requisito determinante: affinché meccanismi di controllo
energetico possano essere integrati senza alterare il modello di
interazione lato utente, l'architettura della piattaforma deve
esporre punti di estensione adeguati a livello di scheduling e
gestione delle funzioni.

In questa prospettiva, la scelta della piattaforma serverless su
cui condurre la sperimentazione non è indifferente. Nella sezione
seguente viene presentata Serverledge, la piattaforma adottata come
riferimento nel presente lavoro, descrivendone le caratteristiche
architetturali rilevanti ai fini della ricerca.

\subsection{La piattaforma Serverledge}
\label{subsec:serverledge}

Serverledge è una piattaforma Function-as-a-Service progettata
per ambienti cloud-edge, caratterizzata da un'architettura
decentralizzata che consente l'esecuzione di funzioni serverless
all'interno di container. La piattaforma separa le componenti di
controllo da quelle di esecuzione, permettendo una gestione
distribuita delle istanze e delle informazioni di stato tra i
nodi che compongono l'infrastruttura.

Ogni nodo Serverledge espone un'interfaccia per la creazione,
l'eliminazione e l'invocazione delle funzioni, supportando sia
la modalità sincrona sia quella asincrona. Le decisioni di
esecuzione sono affidate a uno scheduler, responsabile del
\emph{placement} delle richieste sui nodi disponibili.
L'esecuzione avviene tramite container Docker, mantenuti attivi
per un intervallo di tempo limitato successivamente all'invocazione,
con l'obiettivo di ridurre l'impatto dei \emph{cold start} sulle
invocazioni successive.

La condivisione delle informazioni tra nodi è realizzata tramite
un meccanismo di storage distribuito, che consente di propagare
metadati e informazioni di stato all'interno del cluster. Tali
informazioni sono accessibili dallo scheduler al momento della
decisione di \emph{placement}, permettendo di scegliere tra
l'esecuzione locale sul nodo ricevente e l'offloading verso
altri nodi edge o verso il cloud, in funzione delle condizioni
operative rilevate.

Dal punto di vista architetturale, Serverledge è progettata per
essere estensibile: le logiche di scheduling e di gestione delle
funzioni sono separate dal modello di interazione lato utente,
rendendo possibile l'integrazione di componenti aggiuntivi di
monitoraggio e controllo senza modificare l'interfaccia esposta
agli utenti. Questa proprietà è particolarmente rilevante nel
contesto del presente lavoro, in quanto consente di introdurre
meccanismi di gestione energetica direttamente a livello di
scheduler, sfruttando i punti di estensione nativi della
piattaforma senza alterarne il comportamento di base.